\documentclass[10pt]{article}
\usepackage[margin=0.85in]{geometry}
\usepackage{amsmath,amssymb,amsthm}
\newtheorem{theorem}{Theorem}
\title{A one-edge flag-complex counterexample to conjecture 00000002712}
\author{earthking11}
\date{October 6, 2026}
\begin{document}
\maketitle
\section*{Scope and standard definitions}
The conjecture proposes that the $g$-vectors of flag simplicial complexes
are nonnegative. It does not require the complexes to be homology
spheres, boundaries of polytopes, or closed manifolds. We use a connected,
pure flag complex within that stated class: the full $1$-simplex.
No assertion concerning the usual sphere-restricted $g$-conjecture is
being refuted.

For a $(d-1)$-dimensional complex, let $f_{j-1}$ count faces with $j$
vertices, including $f_{-1}=1$ for the empty face. The standard transform
and truncated $g$-vector are
\[
 \sum_{i=0}^{d}h_i t^{d-i}
 =\sum_{j=0}^{d}f_{j-1}(t-1)^{d-j},\qquad
 g_0=h_0,\quad g_i=h_i-h_{i-1}\quad(1\leq i\leq\lfloor d/2\rfloor).
\]

\begin{theorem}
The clique complex of the actual graph $K_2$ is connected, pure, and
flag, but its standard truncated $g$-vector has $g_1=-1$.
\end{theorem}
\begin{proof}
Let the vertices be $\{0,1\}$ with the single edge $\{0,1\}$. Its clique
complex consists exactly of the empty face, the two singleton faces,
and the full edge. The face collection is downward closed; every set of
pairwise adjacent vertices is a face, so it is flag. The graph is
connected. The full edge is the unique facet, so the complex is pure of
dimension $1$, with rank $d=2$.

Counting these actual faces gives $(f_{-1},f_0,f_1)=(1,2,1)$.
Therefore
\[
 \sum_{i=0}^{2}h_i t^{2-i}
 =(t-1)^2+2(t-1)+1=t^2,
\]
and $h=(1,0,0)$. The truncation range includes index $1$, since
$\lfloor d/2\rfloor=1$. Thus $g_1=h_1-h_0=-1<0$, contradicting the
claimed nonnegativity for flag complexes.
\end{proof}

The counterexample is a simplex (a ball), not a sphere. Adding a sphere
hypothesis would change the original assertion and exclude this example;
the submission does not claim to solve that modified question.
Refuting the stated nonnegativity already suffices, so its further
suggestion of a free-resolution proof is unnecessary.

\section*{Formal verification}
Lean uses Mathlib's actual \texttt{SimpleGraph (Fin 2)} and defines its
clique faces as a finite set of vertex sets. It proves downward closure,
flagness, connectedness, and facet purity. The rank is defined as the
maximum actual face size and proved to equal $2$; the dimension then
equals $1$. All entries of the face counts, including the empty-face
count, are computed from this actual face collection, not assumed as
an external vector.

The $h$-polynomial is an integer polynomial obtained by the displayed
standard transform, and its coefficients define $h_i$. The $g$-vector
has the exact formal index type $0,\ldots,\lfloor d/2\rfloor$.
Lean proves the negative entry and a single combined counterexample
theorem containing the connected, flag, pure, dimension, and negativity
properties. The project compiles on Lean 4.33.1 with pinned Mathlib,
without unfinished proofs, native evaluation, additional axioms, or
warnings. The reported dependencies are propositional extensionality,
classical choice, and quotient soundness. Computation enumerates only
the four faces of this two-vertex graph; no long brute force is used.
\end{document}
